\section{Conclusion}
\label{sec:conclusion}

We have presented Intelligent Coarse Dropout (ICD) and Anti-ICD (AICD), two
augmentation methods that turn a GradCAM saliency map into a tile-based
occlusion mask.  ICD masks the tiles the model finds most salient, removing the
evidence it currently relies on; AICD masks the least salient tiles, perturbing
the background while leaving the discriminative region intact.  The two are
governed by a single percentile threshold that selects them as opposite
directions of the same construction and fixes their occlusion area to a common
fraction.  Masked tiles are replaced by one of several soft fills rather than a
hard box, and the masking operates at a coarse tile granularity that tempers the
noise of pixel-level saliency.  The methods are implemented in the open-source
BNNR library on top of \texttt{pytorch-grad-cam}, with a saliency cache that
makes their online use practical.

The present paper introduces the methods and their rationale; it makes no
empirical claim about their effect on accuracy.  The natural next step is a
controlled comparison of ICD and AICD against RandAugment, TrivialAugment,
AutoAugment, and random masking, under matched compute and a held-out test
split.  That comparison, together with the saliency measurements proposed in
Section~\ref{sec:discussion:saliency}, would test whether the methods reshape the
model's attention as their design intends.  We report that study separately.
Further directions include adapting the tile size to object scale, refreshing
the saliency cache on a schedule, substituting alternative class activation
methods for GradCAM, and extending the tile-scoring construction beyond
classification to detection and segmentation.
